%
%
%
%
%
%
%
%
%
\section*{Use of AI tools}\addcontentsline{toc}{section}{Use of AI tools}

Generative AI tools were used throughout this work: large language models of the Claude family by Anthropic, mainly Claude Opus~5 and~5.5,
and also Claude Fable~5.1 and Claude Sonnet~5.5 and~5, between August and October 2026,
through the Claude Code environment. They were used to draft the text and the mathematical arguments, to write the computer programs, to
run and check the computations, and to search for and read the literature. The author works independently, is not trained as a number
theorist, and learned the subject actively during the work; the tools made it possible to carry out and check computations on this scale. The author led the work in stages: each stage was
set by a prompt that stated what it should achieve, within a stage the tools followed a plan agreed with the author, and the author considered
the results before giving the next prompt, revising the plan when the results required it. The author contributed ideas where they helped and
framed how a problem should be attacked, decided which statements to keep, and is responsible for the content. The programs of the project,
including its own checking and build tools, were developed with the AI tools, and the author started the long computations, such as the
factor searches, on the author's own machine.

Several parts of the checking do not rest on these tools. The elliptic curve primality certificates were produced by PARI/GP~2.17.2 and are
verified by a separate program, and the factor searches use GMP-ECM~7.0.6. Every computed number is generated by a program from deposited
data or from a cited constant, and the programs carry positive and negative controls. Proofs, numbers and citations were also checked in
independent passes by further instances of the AI tools that were given no access to the drafting context. These tools are not authors.
